A soft sensor that reports an uncertainty interval makes two promises, and they are
not the same: that its point estimate is \emph{accurate}, and that its interval is
\emph{valid}, covering the truth at the stated rate. Practice usually conflates them.
We separate the two by construction and then test the separation directly. The
framework pairs physics-derived features and a linear estimator, selected by blocked
time-series cross-validation with a one-standard-error rule, with an open-loop bias
update at the documented laboratory delay and adaptive conformal intervals computed on
the bias-corrected residual; functional scale-bias migration transfers it to new
operating regimes. On two process columns the sensor is both accurate and valid:
held-out $R^2$ rises from $+0.476$ to $+0.857$ on a debutanizer benchmark once drift
is corrected, and on Tennessee Eastman regimes the intervals hold $0.897$--$0.903$
empirical coverage against a $0.90$ target where a static conformal baseline swings
between $0.847$ and $0.957$. The decisive test is a negative control: a deliberately
physics-free process (SECOM, 590 anonymized signals with no thermodynamic structure)
on which the point model is designed to fail. It does ($R^2 = -1.84$), yet conformal
validity is undamaged ($0.910$--$0.915$ coverage at $37\%$ narrower intervals than the
static baseline). Accuracy collapses while validity holds. This dissociation is the
central evidence: distribution-free validity is model-agnostic, whereas point accuracy
is what the physics anchors buy, a claim that can only be established by exhibiting the
case where one survives and the other does not. We further show that under delayed
laboratory feedback the coverage indicator must be scored against the interval issued
when a sample was measured, not the interval the adapted state would issue when the
result arrives, and we enforce this in a real-time implementation that runs two orders
of magnitude inside a typical analyzer cycle.
